\chapter{The \S5.5 Annotation Audit: Draw and Labelling Page}
\label{chap:s2-6}

\section{Goal}

The dataset grades its own annotations, and 47\% of images are self-flagged ``poor info''. Every Study~2
figure is measured against those annotations, so they are a reference standard, not ground truth.
Architecture document \S5.5 asks a person to relabel 100 images, 50 per quality tier, so that the
annotation error rate can be reported beside every result.

\section{What Was Done}

\textbf{The draw.} \texttt{12\_audit\_sample.py} draws from the whole corpus (10,013 images), not the study
sample. The pre-registered draw was a pure random 50 per tier with seed 42. It held 3 hazard images in 100,
too few to estimate how often a \emph{marked} violation is wrong. \textbf{Decided 2026-11-10:} stratify by
hazard within tier --- 25 images with an annotated violation and 25 without, per tier, each pool drawn with
its own seeded generator and independent of record order. Corpus-level error rates are reweighted to the
true 12.8\% hazard prevalence. The 50 hazard images carry 38 rule\_1, 7 rule\_3, 3 rule\_2 and 3 rule\_4
violations, so per-rule rates beyond PPE will rest on very few images. 61 of the 100 images were in the
study sample and carry a model answer.

\textbf{The labelling page.} A private web page presents each image, downscaled to a 1,024-pixel long side.
Step one is blind: the auditor answers the four rule questions (yes / no / can't tell) from the photograph
alone and saves. Only then does the page reveal the dataset's annotation boxes and the model's answer and
region, and the blind answers lock. Step two asks, per annotated box, whether it sits on its hazard (yes /
partly / no), with a free-text note. Each auditor's answers are stored under their own identity; the owner
can read all of them.

\section{Status}

Built and checked empty at owner and contributor access. Awaiting an auditor. A scoring script --- error
rates per tier with intervals, reweighted to prevalence --- is written once answers exist.

\section{Standard vs.\ Non-Standard Choices}

\textbf{Standard.} Blind relabelling of a reference standard; stratified sampling with reweighting.

\textbf{Non-standard, and stated in the paper.} The stratified draw replaces the pre-registered pure random
draw, decided after seeing that the random draw could not answer the question it was for. Blind-then-reveal
in one page keeps the first answers uncontaminated while still collecting box-level judgements.

\section{Implications}

Until the audit is done, every accuracy and explanation figure in Study~2 is ``agreement with the dataset's
annotations''. The audit turns that into an estimate with a known error.
